%! Author = sriadityavedantam
%! Date = 3/25/26

\lesson{39}{Wednesday 12 November 2025 9:10}{Day 39}

\begin{theorem}[Uniform Convergence Theorem]{
    If $f_n \to f_0$ uniformly on $A$, and each $f_n$ is continuous, then $f_0$ is continuous.
}
    Let $c \in A$.

    We want to show
    \[
        \lim_{x \to c} f_0(x) = f_0(c).
    \]

    Let $\eps > 0$ be given.

    Choose $N \in \n$ such that
    \[
        \abs{f_N(x) - f_0(x)} < \frac{\eps}{3} \quad \forall x \in A.
    \]

    Since $f_N$ is continuous at $c$, there exists $\delta > 0$ such that
    \[
        \abs{x - c} < \delta \implies \abs{f_N(x) - f_N(c)} < \frac{\eps}{3}.
    \]

    Then for $\abs{x - c} < \delta$,
    \begin{align*}
        \abs{f_0(x) - f_0(c)}
        &\leq \abs{f_0(x) - f_N(x)} + \abs{f_N(x) - f_N(c)} + \abs{f_N(c) - f_0(c)} \\
        &< \frac{\eps}{3} + \frac{\eps}{3} + \frac{\eps}{3} = \eps.
    \end{align*}
\end{theorem}

\section{Power Series}

\begin{definition}[Power Series]
    A power series centered at $c$ is
    \[
        \sum_{n=0}^\infty a_n (x-c)^n,
    \]
    where $a_0,a_1,\ldots$ are constants.
\end{definition}

\begin{example}
    \begin{gather*}
        \exp(x) \coloneqq \sum_{n=0}^\infty \dfrac{x^n}{n!}, \quad R = \infty.\\
        \sin(x) \coloneqq \sum_{n=0}^\infty (-1)^n \dfrac{x^{2n+1}}{(2n+1)!}.\\
    \end{gather*}
\end{example}

\begin{definition}[Uniformly Cauchy]
    A sequence of functions $(f_n)$ is uniformly Cauchy on $A$ if for every $\eps > 0$, there exists $N \in \n$ such that for all $m,n \geq N$,
    \[
        \abs{f_m(x) - f_n(x)} < \eps \quad \forall x \in A.
    \]
\end{definition}

\begin{proposition}{
    A uniformly Cauchy sequence of functions converges uniformly to some function $g$.
}
    For each $x \in A$, $(f_n(x))$ is a Cauchy sequence in $\r$, hence converges.
    Define
    \[
        g(x) \coloneqq \lim_{n\to\infty} f_n(x).
    \]

    Given $\eps > 0$, choose $N$ such that for all $m,n \geq N$,
    \[
        \abs{f_n(x) - f_m(x)} < \eps \quad \forall x.
    \]

    Fix $n \geq N$ and let $m \to \infty$, then
    \[
        \abs{f_n(x) - g(x)} \leq \eps.
    \]

    So convergence is uniform.
\end{proposition}

\begin{theorem}[Radius of Convergence]{
    Let
    \[
        R \coloneqq \frac{1}{\limsup_{n\to\infty} \sqrt[n]{\abs{a_n}}} \in [0,\infty].
    \]

    Then:
    \begin{itemize}
        \item If $\abs{x-c} < R$, the series $\sum a_n(x-c)^n$ converges absolutely.
        \item If $\abs{x-c} > R$, the series diverges.
        \item If $0 < p < R$, then the series converges uniformly on $[c-p,c+p]$.
    \end{itemize}
}
    Fix $x$.

    By the Root Test,
    \[
        \limsup_{n\to\infty} \sqrt[n]{\abs{a_n(x-c)^n}}
        = \abs{x-c}\cdot \limsup_{n\to\infty} \sqrt[n]{\abs{a_n}}.
    \]

    So convergence occurs when $\abs{x-c} < R$.

    Now let $p < R$.
    Then
    \[
        p \cdot \limsup_{n\to\infty} \sqrt[n]{\abs{a_n}} < 1.
    \]

    So there exists $\gamma \in (0,1)$ and $N$ such that for $n \geq N$,
    \[
        \sqrt[n]{\abs{a_n} p^n} \leq \gamma.
    \]

    Hence
    \[
        \abs{a_n} p^n \leq \gamma^n.
    \]

    For $\abs{x-c} \leq p$,
    \begin{align*}
        \abs{\sum_{k=n}^m a_k (x-c)^k}
        &\leq \sum_{k=n}^m \abs{a_k} p^k \\
        &\leq \sum_{k=n}^m \gamma^k \\
        &= \frac{\gamma^n - \gamma^m}{1-\gamma}.
    \end{align*}

    This goes to $0$ uniformly, so the partial sums are uniformly Cauchy.

    Hence the series converges uniformly on $[c-p,c+p]$.
\end{theorem}